
\chapter{Intermission and review}

Welcome to Volume Two of the Crystal Ball series: \textit{Foundations for Data
Science}! I titled the first volume ``\textit{Introduction to Data Science}''
because it led you through a dip-your-toes-in-the-water experience. You took a
brief tour through the various elements in this diverse field and got a feel
for what it was all about.

Since you're still reading, this means you're still interested, and ready to
explore the next level. In \textit{Foundations}, we'll solidify your growing
knowledge so that you have a firm base on which to build everything else.
Future volumes in this series will cover applications, advanced techniques, and
some special data types that require special treatment. But it all ultimately
rests on the skill set you'll have built in Volumes One and Two of this series.

We're going to dive right in. But first, let me make a list of some of the
things I'm counting on you being pretty solid on:

\begin{compactitem}
\small
\itemsep.1em
\item Atomic and aggregate variables

\begin{compactitem}
\item \texttt{int}s, \texttt{float}s, \texttt{strings}
\item NumPy \texttt{array}s, Pandas \texttt{Series}es, \texttt{DataFrame}s
\item What all these look like in memory
\end{compactitem}

\item Scales of measure (especially categorical vs.~numeric)

\item Assocation vs.~causality

\begin{compactitem}
\item Causal diagrams
\item Dependent and independent variables
\item Observational studies vs.~controlled experiments
\item Confounding factors
\end{compactitem}

\item Statistical significance

\item Exploratory Data Analysis
\vspace{-.18in}
\begin{multicols}{2}
    \begin{compactitem}
    \item Bar charts
    \item Histograms
    \item Contingency tables
    \item Scatterplots
    \item Boxplots
    \item Quantiles
    \item Outliers
    \end{compactitem}
\end{multicols}
\vspace{-.18in}

\item Python stuff
    \begin{compactitem}
    \item Creating and performing calculations on atomic data
    \item Creating and accessing NumPy arrays
    \item \texttt{for} loops
    \item \texttt{if} statements (and \texttt{if}/\texttt{elif}/\texttt{else})
    \item Functions and methods
        \begin{compactitem}
        \item Calling a function/method vs.~writing a function
        \item Passing arguments (both atomic and aggregate)
        \item Return values
        \item Function that modify in-place vs.~returning a copy
        \end{compactitem}

    \item Pandas \texttt{Series} \& \texttt{DataFrames}
        \begin{compactitem}
        \itemsep.1em
        \item Reading from a \texttt{.csv} file
        \item The ``index''
        \item \texttt{.value\_counts()}
        \item \texttt{.groupby()}
        \item Using \texttt{.iloc[]} vs.~\texttt{.loc[]} vs. \texttt{[]}
        \item Single ints/labels vs.~slices vs.~lists
        \item Queries
        \item Recoding / transforming columns
        \end{compactitem}
    \end{compactitem}

\item ML concepts
    \begin{compactitem}
    \item Classification vs.~regression
    \item Features and target attribute
    \item Training data, test data, and new data
    \item Random sampling
    \item The ``prior'' vs.~``posterior''
    \item Evaluating a classifier
    \end{compactitem}

\end{compactitem}

All these topics were covered in \textit{Crystal Ball} Volume One. If any of
them are iffy, you might take an hour or two to flip back through them and
brush up!
